\documentclass[hyperref={pdfpagelabels=false, unicode},pdf,slideColor,fyma,9pt]{beamer}
\usepackage[utf8]{inputenc}		%kodovani zdrojoveho textu
\usepackage[english]{babel} 		%cestina
\usepackage[IL2]{fontenc}		%font vyladeny pro cestinu
\usepackage{lmodern}			%odstraneni varovani pro IL2
\usepackage{graphicx}
\usepackage{color}
\usepackage{multirow}
\usepackage{subfigure}

% \usepackage[UKenglish]{isodate}

\newcommand\myqt[1]{\quotedblbase #1\textquotedblleft}

%cislovani stran:



\mode<presentation> {
  \usetheme{Boadilla}
	\usecolortheme{default}

  %\setbeamercovered{Berlin beaver default}
	%\setbeamertemplate{navigation symbols}{} % To remove the navigation symbols from the bottom of all slides uncomment this line
%\setbeamertemplate{headline}{}
}

\makeatother
\setbeamertemplate{footline}
{
  \leavevmode%
  \hbox{%
  \begin{beamercolorbox}[wd=.3\paperwidth,ht=2.25ex,dp=1ex,center]{author in head/foot}%
    \usebeamerfont{author in head/foot}\insertshortauthor
  \end{beamercolorbox}%
  \begin{beamercolorbox}[wd=.6\paperwidth,ht=2.25ex,dp=1ex,center]{title in head/foot}%
    \usebeamerfont{title in head/foot}\insertshorttitle\hspace*{3em}
  \end{beamercolorbox}}%
  \begin{beamercolorbox}[wd=.1\paperwidth,ht=2.25ex,dp=1ex,center]{date in head/foot}%
    \insertframenumber{} / \inserttotalframenumber\hspace*{1ex}
  \end{beamercolorbox}%
  \vskip0pt%
}
\makeatletter
\setbeamertemplate{navigation symbols}{}


\title{Performance Testing and Analysis of Qpid-Dispatch Router}
% \subtitle{Master Thesis}
\author[Bc. Jakub Stejskal]{Bc. Jakub Stejskal}
\institute[]{Faculty of Information Technology}
\date[]{\today}

\begin{document}
		\begin{frame}
				\titlepage
        \begin{figure}[ht]
        	\begin{center}
        		\scalebox{0.25}{\includegraphics{figures/maestro-logo.pdf}}
        	\end{center}
        \end{figure}
		\end{frame}

		%/////////////////////////////////////////////////
		\section{Motivation}
		\begin{frame}{Motivation}
      \begin{tabular}{cl}
        \begin{tabular}{l}
            \parbox{0.45\linewidth}{%  change the parbox width as appropiate
            \begin{itemize}
            \setlength\itemsep{0.5em}
              \item What is Maestro?
              \vspace{0.5em}
              \begin{itemize}
                  \setlength\itemsep{0.5em}
                  \item Performance tool for \textbf{Message-oriented middleware}.
                  \item Automated, provides reporting.
                  \item Single node, multi node or cluster testing.
              \end{itemize}
              \vspace{0.5em}
              \item Capabilities
              \vspace{0.5em}
              \begin{itemize}
                  \setlength\itemsep{0.5em}
                  \item \textbf{Throughput} measurements.
                  \item \textbf{Latency} measurements.
              \end{itemize}
            \end{itemize}
            }
        \end{tabular}
        & \begin{tabular}{l}
            \scalebox{0.25}{\includegraphics{figures/msg_perf_tool.pdf}}
          \end{tabular} \\
      \end{tabular}
		\end{frame}
		%/////////////////////////////////////////////////
		\section{Extension: Maestro-Agent and AMQP Inspector}
		\begin{frame}{Extensions: Maestro-Agent and AMQP Inspector}
      \begin{tabular}{cl}
        \begin{tabular}{l}
            \parbox{0.55\linewidth}{%  change the parbox width as appropiate
            \begin{itemize}
                \setlength\itemsep{0.5em}
                \item Enables automatic \textbf{SUT} (Software Under Test) changes during the test.
                \item \textbf{Groovy} code handler.
                \item Can fetch and process external repositories.
                \item Monitors \textbf{Qpid-dispatch}.
                \item Request-response mechanism over \textbf{AMQP}.
                \item \textbf{Topology Generator} for automatic deployment and configuration.
            \end{itemize}
            }
        \end{tabular}
        & \begin{tabular}{l}
            \scalebox{0.40}{\includegraphics{figures/agent_demo.pdf}}
          \end{tabular} \\
      \end{tabular}
		\end{frame}
		%/////////////////////////////////////////////////
		% \section{Extension: Topology Generator}
		% \begin{frame}{Extension: Topology Generator}
    %    		\begin{tabular}{cl}
    %        		\begin{tabular}{l}
    %          		\parbox{0.4\linewidth}{%  change the parbox width as appropiate
		% 				\begin{itemize}
		% 					\setlength\itemsep{0.5em}
		% 					\item Metadata transformed into Configuration files.
    %           \begin{itemize}
    %             \item Inventory
    %             \item Graph File
    %           \end{itemize}
		% 					\item Automatic deployment by \textbf{Ansible}.
		% 				\end{itemize}
   % 					}
    %      		\end{tabular}
    %      		& \begin{tabular}{l}
    %        			\scalebox{0.3}{\includegraphics{figures/generator.pdf}} \\ \\ \\
    %        			\scalebox{0.3}{\includegraphics{figures/deployment.pdf}}
    %        		\end{tabular} \\
		% 	\end{tabular}
		% \end{frame}
    %/////////////////////////////////////////////////
    \section{Measurements}
    \begin{frame}{Experimental Evaluation}
      \begin{itemize}
          \setlength\itemsep{0.5em}
          \item \textbf{Latency} vs. \textbf{Throughput} vs. \textbf{Behavioral} testing.
          \item The figure shows throughput comparison of \textbf{Messaging Broker} and \textbf{Qpid-dispatch}.
      \end{itemize}
      \begin{figure}[ht]
        \begin{center}
          \scalebox{0.7}{\includegraphics{figures/throughput.pdf}}
        \end{center}
      \end{figure}
    \end{frame}
    %/////////////////////////////////////////////////
		\section{Measurements II.}
    \begin{frame}{Experimental Evaluation}
      \begin{itemize}
          \setlength\itemsep{0.5em}
          \item Latency is measured by \textbf{Maestro Receiver}.
          \item The figure shows latency comparison between \textbf{Messaging Broker} and \textbf{Qpid-dispatch}.
      \end{itemize}
      \begin{figure}[ht]
        \begin{center}
          \scalebox{0.7}{\includegraphics{figures/latency.pdf}}
          % \scalebox{0.4}{\includegraphics{figures/agent.pdf}}
        \end{center}
      \end{figure}
    \end{frame}
    %/////////////////////////////////////////////////
    \section{Measurements III.}
    \begin{frame}{Experimental Evaluation}
      \begin{itemize}
          \setlength\itemsep{0.5em}
          \item Latency and throughput can be affected by the \textbf{Agent}.
          \item The figure shows behavioral test case, during which one node of topology is shut down.
      \end{itemize}
      \begin{figure}[ht]
        \begin{center}
          \scalebox{0.7}{\includegraphics{figures/agent_latency.pdf}}
        \end{center}
      \end{figure}
    \end{frame}
		%/////////////////////////////////////////////////
		\section{Summary}
		\begin{frame}{Summary}
				\begin{itemize}
						\setlength\itemsep{0.5em}
						\item Extensions of the Maestro
						\vspace{0.5em}
						\begin{itemize}
							\setlength\itemsep{0.5em}
							\item \textbf{Maestro-Agent}
							\item \textbf{AMQP Inspector}
              \item \textbf{Topology Generator} (external tool)
						\end{itemize}
						\vspace{1em}
						\item Experimental Evaluation
						\vspace{0.5em}
						\begin{itemize}
							\setlength\itemsep{0.5em}
							\item \textbf{Throughput} and \textbf{latency}, \textbf{behavioral} testing.
              \item Single node, multi node or cluster topologies.
              \item Topologies consists of \textbf{Qpid-dispatch} and other messaging services.
              \item Automatic generation of specific topologies.
						\end{itemize}
				\end{itemize}
		\end{frame}
		%/////////////////////////////////////////////////
    \section{Contacts}
    \begin{frame}{Contacts}
      \begin{figure}[ht]
        \begin{center}
          \scalebox{0.30}{\includegraphics{figures/maestro-logo.pdf}}
        \end{center}
      \end{figure}
      \begin{center}
        \begin{itemize}
          \item Jakub Stejskal\,---\,jstejska@redhat.com, xstejs24@stud.fit.vutbr.cz
          \item Otavio Rudolfo Piske\,---\,opiske@redhat.com (Author of Maestro)
          \item GitHub\,---\,\url{https://github.com/maestro-performance/}
        \end{itemize}
    \end{center}
    \end{frame}
    %/////////////////////////////////////////////////
\end{document}
